\section{The Costs of Throwing an Exception}\label{sec:ExceptionCosts}
To understand this section, some background knowledge regarding the interaction between the Java compiler \texttt{javac}\index{javac}, the \Java Virtual Machine (JVM), and the integrated JIT\index{JIT} compiler (HotSpot\index{HotSpot} for Oracle's JVM) is helpful.

The Java compiler \texttt{javac}\index{javac} translates Java source code into Java bytecode ­– a form of \bdq{machine language} executed interpretively within a virtual machine (the JVM) ­– entailing the performance issues inherent to interpreters.

As mentioned, the JIT\index{JIT} (just-in-time) compiler is a component of the JVM; it translates Java bytecode into native machine code at runtime, while the program is executed. However, only code that is executed repeatedly is translated; furthermore, the generated machine code may be discarded and replaced with newly generated code as part of further optimizations.

This mechanism is responsible for the steady increase in a Java program's throughput over time until a saturation point is reached.

It should be equally obvious that throughput collapses when the JVM must discard the generated machine code and revert to executing the Java bytecode in interpreter mode.

%------------------------------------------------------------------------------

\subsection{What happens when an Exception is thrown?}
\begin{enumerate}
	\item{Object Creation and Stack Trace Capture
	 
	When \bdq{\lstinline|throw new FooException()|} is executed, normal object allocation happens, then the constructor\index{constructor} of \lstinline|java.lang.Throwable|\index{java.lang!Throwable} \autocite{ORACLE_DOC_THROWABLE_CLASS} calls the native method \lstinline|fillInStackTrace()|\index{java.lang!Throwable!fillInStackTrace()} \autocite{ORACLE_DOC_THROWABLE:fillInStackTrace}. This runs synchronously in the throwing thread itself (no safepoint is needed, since it is inspecting its own stack) and builds a compact internal backtrace structure – not yet the finished \lstinline|StackTraceElement[]|\index{java.lang!StackTraceElement} array \autocite{ORACLE_DOC_STACKTRACEELEMENT_CLASS}. That array only materializes lazily when \lstinline|Throwable::getStackTrace|\index{java.lang!Throwable!getStackTrace()} \autocite{ORACLE_DOC_THROWABLE:getStackTrace} or \lstinline|Throwable::printStackTrace|\index{java.lang!Throwable!printStackTrace()} \autocite{ORACLE_DOC_THROWABLE:printStackTrace} is called.
	
	For the JVM-internal exceptions
	\begin{itemize}[nosep]
		\item{\lstinline|java.lang.ArithmeticException|\index{java.lang!ArithmeticException} \autocite{ORACLE_DOC_ARITHMETICEXCEPTION_CLASS}}
		\item{\lstinline|java.lang.ArrayIndexOutOfBoundsException|\index{java.lang!ArrayIndexOutOfBoundException} \autocite{ORACLE_DOC_ARRAYINDEXOUTOFBOUNDSEXCEPTION_CLASS}}
		\item{\lstinline|java.lang.ArrayStoreException|\index{java.lang!ArrayStoreException} \autocite{ORACLE_DOC_ARRAYSTOREEXCEPTION_CLASS}}
		\item{\lstinline|java.lang.ClassCastException|\index{java.lang!ClassCastException} \autocite{ORACLE_DOC_CLASSCASTEXCEPTION_CLASS}}
		\item{\lstinline|java.lang.NullPointerException|\index{java.lang!NullPointerException} \autocite{ORACLE_DOC_NULLPOINTEREXCEPTION_CLASS}}
	\end{itemize}
	there is a specific \bdq{optimization} JVM parameter that is on by default:
	\begin{center}
		\texttt{\mbox{-XX:+OmitStackTraceInFastThrow}}\index{OmitStackTraceInFastThrow}  
	\end{center}
	If the same implicit Exception\index{exception} is thrown repeatedly (\bdq{hot}) at the same bytecode location, HotSpot\index{HotSpot} replaces the allocation and the stack-trace-fill step with a preallocated, stack-trace-less instance (\bdq{\texttt{… <no stack trace available>}}). This is the classic reason exceptions\index{exception} used as control flow in loops eventually show up with no stack trace at all – and it is also why \lstinline|NullPointerException|'s helpful detailed messages disappear once this kicks in, since the backtrace data the message relies on is missing.
	
	Benchmarks show a call to \lstinline|new Exception()|\index{java.lang!Exception} to instantiate an Exception\index{exception} object costs roughly 1200ns versus around 8ns for a simple object like \lstinline|java.util.Date|\index{java.util!Date} (and \lstinline|Date| is far from being really simple).
	
	The calls to \lstinline|Throwable::fillInStackTrace|\index{java.lang!Throwable!fillInStackTrace()} causes a measurable CPU load under a high exception rate (e.g., when validation exceptions\index{exception} are used to control the execution flow). In addition, the costs for filling the stack trace scales with stack depth that is often in the range of \mbox{50+} frames for code executed within an application server or web container. Of course this affects also the next step.}
	
	\item{\lstinline|athrow|\index{athrow} and Stack Unwinding 
	
	The \lstinline|athrow|\index{athrow} bytecode instruction pops the reference from the stack and starts the handler\footnote{In source code, that \bdq{handler} is a \bsq{\lstinline|catch( final <Exception> e )|}\index{try!catch} line in most cases, where \texttt{<Exception>} matches the type of the thrown Exception\index{exception}, or a parent class.} search: the JVM (interpreter or compiled code, via the method's exception table/OopMaps\index{OopMaps}) checks whether the current frame has a matching handler for the thrown exception\index{exception}. The type check is handled via the exception’s\index{exception} class hierarchy.
	
	If no match could be found, the frame is popped from the stack, during which:
	\begin{itemize}
		\item{\lstinline|monitorexit|\index{monitorexit} (another bytecode instruction) runs for every active synchronized\index{synchronized} block/method\index{method!synchronized} in that frame, to release the locks}
		\item{Local variables become eligible for the garbage collection}
		\item{The search continues recursively in the caller frame}
	\end{itemize}
	If no matching handler exists anywhere on the stack, the internal method \lstinline|Thread::dispatchUncaughtException|\index{java.lang!Thread!dispatchUncaughtException()} will call the implementation of the interface \lstinline|UncaughtExceptionHandler|\index{java.lang!Thread.UncaughtExceptionHandler} \autocite{ORACLE_DOC_UNCAUGHTEXCEPTIONHANDLER_INTERFACE} that is either set to the \lstinline|Thread|\index{java.lang!Thread} \autocite{ORACLE_DOC_THREAD_CLASS} object itself or to the associated \lstinline|ThreadGroup|\index{java.lang!ThreadGroup} \autocite{ORACLE_DOC_THREADGROUP_CLASS}. If none is set, the parent thread is searched for an \lstinline|UncaughtExceptionHandler|\index{java.lang!Thread.UncaughtExceptionHandler}.}
\end{enumerate}

%------------------------------------------------------------------------------

\subsection{JIT Deoptimization}\label{sec:JITDeopt}
This is where the real cost lies if things go wrong.
 
The JIT\index{JIT} compiler profiles exception paths. If a \lstinline|catch|\index{try!catch} branch is rarely or never taken during profiling, the compiler often doesn't compile it in fully – it places an \textit{uncommon} trap there instead. If that path is then actually hit at runtime, it triggers a \textit{deoptimization of that specific method}\index{method}: execution falls back to the interpreter for that frame and continues interpreted from there.

\bdq{Deoptimization} means the compiler's optimistic assumption was proven wrong, so execution switches from compiled code back to the interpreter by reconstructing interpreter state, and the compiled code may either be discarded or recompiled.

Whether the affected method\index{method} gets recompiled depends on the trap history (\texttt{PerMethodTrapLimit}/\texttt{PerBytecodeTrapLimit}). If the trap fires too often, HotSpot\index{HotSpot} permanently marks the method as no longer compilable for that path at full optimization – it stays interpreted or degraded.

Simply throwing/catching exceptions through code that was already compiled with the handler correctly baked in from the start does \textit{not} cause deoptimization – unwinding through compiled frames runs through generated runtime stubs, GC-safe via OopMaps, without falling back to the interpreter.

In short: deoptimization doesn't happen because of the \lstinline|throw| itself, but when the exception\index{exception} activates an execution path the compiler had classified as \bdq{cold} at compile time.

Practical takeaway: exceptions\index{exception} for genuine error cases (rare, truly exceptional) are essentially free in practice. Exceptions\index{exception} used as regular control flow (e.g., validation on every request) are the case where stack trace cost, garbage collection pressure, and deoptimization churn actually matter – there it's often worth avoiding exceptions\index{exception} as a signaling mechanism.

%==============================================================================
% End of file!!
%==============================================================================

